\section{The passage is the unit of comparison}
\label{sec:scope}

I use \emph{situated} to describe prose shaped for a particular reader and
purpose. Its author anticipates likely objections, selects facts according to
what that reader already knows, and preserves uncertainty that the evidence
cannot resolve. Language agents can imitate every one of these choices. When a
prompt leaves the situation underspecified, however, their answers can
converge on a portable middle register that could address many readers and
fully satisfy none of them.

The pattern catalog is descriptive: it records recurrent editorial problems
but does not estimate their frequency. Any comparison is therefore
probabilistic and passage-level. Grammar, formal vocabulary, and conventional
punctuation provide weak evidence about provenance because each has ordinary
explanations. A clean sentence establishes little about its author.

\subsection{Audience specification changes the choice of detail}

The intended reader determines which background matters and where a comparison
becomes useful. Experimental controls matter differently to a reviewer
checking the evidence than to a student learning the mechanism. A collaborator
choosing between methods needs the governing tradeoff earlier. The technical
subject may be identical, but the information problem changes with its use.

Generic management language obscures this dependence. ``Give each section one
job'' works across many contexts because it does not identify a reading
situation.
``Each section should help a reviewer answer one question'' is narrower: it
connects the section's organization to a particular act of evaluation. The
review task determines the structure; merely adding the word \emph{reader}
would not make the sentence specific.

\subsection{Pattern claims remain properties of passages}

Stylistic labels should describe properties that can be checked against the
passage. A statement about the person or system that produced it requires
independent evidence. This distinction matters because style-based detection
has familiar false positives: professional editing explains consistent
grammar, word processors explain curly quotation marks, and a field's shared
vocabulary explains many repeated terms.

Clusters still matter, but for diagnosis rather than attribution. Inflated
emphasis becomes more damaging when it coincides with generic attribution and
a paragraph whose sentences all share one shape. Together, those choices can
hide the claim and its support. They remain evidence about the prose, not proof
of authorship.
